\Article{15}
In der vorhergehenden Abhandlung, Art. 37 und folgende, haben wir eine Methode angegeben, mit der sich die Genauigkeit der Beobachtungen nachträglich möglichst genau bestimmen lässt. Sind nämlich die Näherungswerte von \(\pi\) Größen durch Beobachtungen gleicher Genauigkeit bekannt geworden und werden sie mit denjenigen Werten verglichen, die durch Rechnung aus den plausibelsten Werten der \(\rho\) Elemente hervorgehen, von denen jene Größen abhängen, so sind die Quadrate der Unterschiede zu addieren und ihre Summe durch \(\pi-\rho\) zu dividieren. Der so erhaltene Quotient kann dann als Näherungswert des Quadrats des dieser Art von Beobachtungen anhaftenden mittleren Fehlers angesehen werden. Haben die Beobachtungen ungleiche Genauigkeit, so sind diese Regeln nur insofern abzuändern, als die Quadrate vor der Addition mit den Gewichten der Beobachtungen multipliziert werden müssen und der auf diese Weise erhaltene mittlere Fehler auf Beobachtungen zu beziehen ist, deren Gewicht als Einheit angenommen wurde.

In der vorliegenden Behandlung stimmt jene Summe offenbar mit der Summe \(S\) überein und der Unterschied \(\pi-\rho\) mit der Anzahl der Bedingungsgleichungen \(\sigma\). Daher erhalten wir für den mittleren Fehler von Beobachtungen, deren Gewicht \(=1\) ist, den Ausdruck \(\BareRoot{\frac S\sigma}\); diese Bestimmung verdient umso größeres Vertrauen, je größer die Anzahl \(\sigma\) ist.

Es wird jedoch der Mühe wert sein, dies auch unabhängig von der vorhergehenden Untersuchung festzustellen. Dazu wird es zweckmäßig sein, einige neue Bezeichnungen einzuführen. Es sollen nämlich den folgenden Werten der Unbestimmten \(\xi,\eta,\zeta\) etc.,
\SourceLeaf{079}
\[
\xi=a,\quad\eta=b,\quad\zeta=c\quad\text{etc.}
\]
diese Werte von \(x,y,z\) etc. entsprechen,
\[
x=\alpha,\quad y=\beta,\quad z=\gamma\quad\text{etc.}
\]
sodass
\[
\begin{aligned}
\alpha&=a[\alpha\alpha]+b[\alpha\beta]+c[\alpha\gamma]+\text{etc.}\\
\beta&=a[\alpha\beta]+b[\beta\beta]+c[\beta\gamma]+\text{etc.}\\
\gamma&=a[\alpha\gamma]+b[\beta\gamma]+c[\gamma\gamma]+\text{etc.}
\end{aligned}
\]
etc. gilt. Ebenso nehmen wir an, dass den Werten
\[
\xi=a',\quad\eta=b',\quad\zeta=c'\quad\text{etc.}
\]
diese entsprechen,
\[
x=\alpha',\quad y=\beta',\quad z=\gamma'\quad\text{etc.}
\]
und diesen,
\[
\xi=a'',\quad\eta=b'',\quad\zeta=c''\quad\text{etc.}
\]
die folgenden,
\[
x=\alpha'',\quad y=\beta'',\quad z=\gamma''\quad\text{etc.}
\]
und so weiter.

Unter diesen Voraussetzungen liefert die Verbindung der Gleichungen (4) und (9)
\[
\begin{aligned}
A&=\alpha e+\alpha'e'+\alpha''e''+\text{etc.}\\
B&=\beta e+\beta'e'+\beta''e''+\text{etc.}\\
C&=\gamma e+\gamma'e'+\gamma''e''+\text{etc.}
\end{aligned}
\]
etc. Da somit \(S=\mathfrak A A+\mathfrak B B+\mathfrak C C+\text{etc.}\) gilt, ergibt sich offenbar
\[
\begin{aligned}
S={}&(ae+a'e'+a''e''+\text{etc.})(\alpha e+\alpha'e'+\alpha''e''+\text{etc.})\\
&+(be+b'e'+b''e''+\text{etc.})(\beta e+\beta'e'+\beta''e''+\text{etc.})\\
&+(ce+c'e'+c''e''+\text{etc.})(\gamma e+\gamma'e'+\gamma''e''+\text{etc.})+\text{etc.}
\end{aligned}
\]
\Article{16}
Die Durchführung der Beobachtungen, durch die wir Werte der Größen \(v,v',v''\) etc. erhalten, die mit zufälligen Fehlern \(e,e',e''\) etc. behaftet sind, können wir als
\SourceLeaf{080}
ein Experiment betrachten, das zwar die Größe jedes einzelnen begangenen Fehlers nicht erkennen lässt, aber bei Anwendung der oben erläuterten Regeln den Wert der Größe \(S\) liefert, der nach der eben gefundenen Formel eine gegebene Funktion jener Fehler ist. Bei einem solchen Experiment können die zufälligen Fehler natürlich teils größer, teils kleiner ausfallen; je mehr Fehler jedoch zusammentreffen, desto größer ist die Hoffnung, dass der Wert der Größe \(S\) in einem einzelnen Experiment nur wenig von seinem Mittelwert abweicht. Es kommt also wesentlich darauf an, den Mittelwert der Größe \(S\) festzustellen. Nach den in der vorhergehenden Abhandlung dargestellten Grundsätzen, deren Wiederholung hier überflüssig wäre, finden wir diesen Mittelwert
\[
\begin{aligned}
&=(a\alpha+b\beta+c\gamma+\text{etc.})mm\\&\qquad+(a'\alpha'+b'\beta'+c'\gamma'+\text{etc.})m'm'\\
&\qquad+(a''\alpha''+b''\beta''+c''\gamma''+\text{etc.})m''m''+\text{etc.}
\end{aligned}
\]
\Needspace{20\baselineskip}

Bezeichnen wir den mittleren Fehler von Beobachtungen, deren Gewicht \(=1\) ist, mit \(\mu\), sodass \(\mu\mu=pmm=p'm'm'=p''m''m''\) etc. gilt, so lässt sich der eben gefundene Ausdruck wie folgt darstellen:
\[
\begin{aligned}
&\left(\frac{a\alpha}{p}+\frac{a'\alpha'}{p'}+\frac{a''\alpha''}{p''}\ \text{etc.}\right)\mu\mu
\\&\qquad+\left(\frac{b\beta}{p}+\frac{b'\beta'}{p'}+\frac{b''\beta''}{p''}+\text{etc.}\right)\mu\mu\\
&\qquad+\left(\frac{c\gamma}{p}+\frac{c'\gamma'}{p'}+\frac{c''\gamma''}{p''}+\text{etc.}\right)\mu\mu+\text{etc.}
\end{aligned}
\]
Die Summe \(\frac{a\alpha}{p}+\frac{a'\alpha'}{p'}+\frac{a''\alpha''}{p''}+\text{etc.}\) ergibt jedoch
\[
=[aa].[\alpha\alpha]+[ab].[\alpha\beta]+[ac].[\alpha\gamma]+\text{etc.}
\]
und ist daher \(=1\), wie aus dem Zusammenhang der Gleichungen (6) und (7) leicht ersichtlich ist. Ebenso ist
\[
\begin{aligned}
\frac{b\beta}{p}+\frac{b'\beta'}{p'}+\frac{b''\beta''}{p''}+\text{etc.}&=1\\[5pt]
\frac{c\gamma}{p}+\frac{c'\gamma'}{p'}+\frac{c''\gamma''}{p''}+\text{etc.}&=1
\end{aligned}
\]
und so weiter.

Damit wird schließlich der Mittelwert von \(S\) \(=\sigma\mu\mu\); soweit also der zufällige Wert von \(S\) als sein Mittelwert angenommen werden darf, ist \(\mu=\BareRoot{\frac S\sigma}\).

\clearpage
